% ------------------------------------------------
\StartSection{產生論文 Generate Thesis}{chapter:how-to:use:generate}
% ------------------------------------------------

把LaTeX變成PDF有三種方式, 選一種即可. 三種方式產生的PDF相同, 差別只在於要不要安裝軟體, 以及習慣用什麼來編輯.

\begin{DescriptionFrame}
\begin{verbatim}
方式一  Overleaf         不用安裝軟體, 在瀏覽器內編寫和產生
方式二  本機安裝+終端機   TeX Live/MacTeX/MiKTeX, 一行指令產生
方式三  本機安裝+編輯器   同方式二安裝, 再用Texmaker/TeXstudio產生
\end{verbatim}
\end{DescriptionFrame}

不論哪一種方式, 正式產生論文的指令都是同一句. \texttt{latexmk}會自動執行XeLaTeX和BibTeX, 並重複到目錄, 引用和參考文獻都正確為止, 不需要自己猜要跑幾次:

\begin{center}
\verb|latexmk -xelatex -synctex=1 -interaction=nonstopmode thesis.tex|
\end{center}

% ------------------------------------------------
\StartSubSection{方式一: Overleaf}

Overleaf是線上的LaTeX環境, 本模版在Overleaf Gallery有公開範本, 連結放在本專案GitHub頁面\RefBib{web:this-project:github}的最上方.

\begin{enumerate}
  \item 開啟範本頁面, 按`Open as Template', 範本會複製到你的Overleaf帳號.
  \item 在左上角Menu確認Compiler是XeLaTeX, Main document是`thesis.tex'.
  \item 按Recompile. 要完整重建時, 在Recompile旁的選單選`Recompile from scratch'.
\end{enumerate}

Overleaf上的範本已經關閉`\verb|\ExampleMode|'和初稿標記, 直接填寫`conf/conf.tex'即可. 免費方案有編譯時間限制, 所以不要在Overleaf開啟`\verb|\ExampleMode|'去產生這份教學文件. Overleaf範本和GitHub Releases是分開更新的, Overleaf上的版本可能較舊; 需要最新修正時, 請下載GitHub Releases的學生套件在本機產生.

% ------------------------------------------------
\StartSubSection{方式二: 本機安裝TeX}

本機需要一個包含XeLaTeX, BibTeX和\texttt{latexmk}的TeX發行版, 以下三個都可以:

\begin{description}
  \item[TeX Live (Windows / Linux)] \hfill\\
    到TeX Live網頁 (\url{https://www.tug.org/texlive/}) 下載安裝程式, 選擇完整安裝(full scheme), 就不會缺少套件. 完整安裝需要數GB的空間和一段時間.

  \item[MacTeX (macOS)] \hfill\\
    到MacTeX網頁 (\url{https://www.tug.org/mactex/}) 下載安裝, 內容就是macOS版的TeX Live完整安裝.

  \item[MiKTeX (Windows)] \hfill\\
    到MiKTeX網頁\RefBib{web:miktex:website}下載安裝. 安裝時把`Install missing packages on-the-fly'設為Yes, 第一次產生論文時MiKTeX會自動下載需要的套件, 所以會等比較久. 之後定期在MiKTeX Console按Update更新套件; 看到某個package太舊的錯誤訊息時, 就是要更新了.
\end{description}

安裝完後開一個終端機(Windows的命令提示字元或PowerShell, macOS的Terminal, Linux的任何終端機), 輸入\verb|latexmk --version|和\verb|xelatex --version|, 兩個都有顯示版本就代表安裝成功.

要產生論文時, 在終端機切換到包含`thesis.tex'的資料夾, 執行上面那句\texttt{latexmk}指令, 就會產生`thesis.pdf'. 電子版只需要這個檔案, 它已包含封面內頁. 印刷版另外需要交給影印店的封面檔, 把指令中的`thesis.tex'換成`cover.tex'即可, 會產生`cover.pdf'.

% ------------------------------------------------
\StartSubSection{方式三: 編輯器}

方式三同樣要先完成方式二的安裝, 之後在編輯器內按一個按鈕就能產生PDF並預覽.

\begin{description}
  \item[Texmaker / TeXstudio] \hfill\\
    在Options (選項) -> Configure -> Commands, 把其中一個編譯指令(如LaTeX或Build)設為\\
    \verb|latexmk -xelatex -synctex=1 -interaction=nonstopmode %.tex|\\
    並把`thesis.tex'設為主文件(Master / Root document). 之後不論正在編輯哪個檔案, 按編譯都會產生整份論文.

  \item[VS Code + LaTeX Workshop] \hfill\\
    安裝LaTeX Workshop擴充功能, 在設定中新增一個recipe, 指令為\\
    \verb|latexmk -xelatex -synctex=1 -interaction=nonstopmode %DOC%|\\
    並以`thesis.tex'為root file.
\end{description}

編輯檔案時請用支援UTF-8的編輯器, 上述的編輯器都是. 用Windows內建記事本的舊版本儲存會變成ANSI編碼, 中文就會變成亂碼.

% ------------------------------------------------
\StartSubSection{產生PDF的流程}

在編譯LaTeX成PDF時, 必須讓目錄, 交叉引用及參考文獻(例如\verb|\RefBib{}|)完成所需的多輪處理. \texttt{latexmk}會讀取dependency及auxiliary files, 只執行當次需要的XeLaTeX及BibTeX passes; 一般正文修改通常不需要重新執行BibTeX.

寫作時可以保持continuous preview, 儲存任何被引用的TeX, BibTeX或圖片檔後, PDF會自動重建:

\begin{center}
\verb|latexmk -xelatex -pvc -view=none -synctex=1 thesis.tex|
\end{center}

\verb|-view=none|避免另外打開第二個viewer; 可繼續使用編輯器內置, 能自動重新載入的PDF viewer. 停止監察可按\texttt{Ctrl-C}.

實際編寫自己論文時, 請在\verb|conf/conf.tex|關閉\verb|\ExampleMode|.

主文件會改用\verb|context/context.tex|; 完整教學範例包含大量圖片及附錄PDF, 重建時間會較長.

如果auxiliary files損壞, 可執行\verb|latexmk -C thesis.tex|清除可重新產生的檔案, 再重新build. 不要把反覆刪檔及猜測XeLaTeX/BibTeX次序當成正常流程. 正式提交前應再執行一次普通\texttt{latexmk} build, 並檢查完整PDF, 頁碼, 目錄, 引用及參考文獻.

% ------------------------------------------------
\StartSubSection{遇到錯誤時}

看log時先找第一個以`!'開頭的行, 那才是真正的錯誤; 後面的訊息多數只是連鎖反應. 常見的錯誤訊息, 原因和解法整理在學生套件的`README.md'中`常見錯誤'一節, 請先在那邊找.

% ------------------------------------------------
